\section{Conclusion}\label{s:conclusion}

This thesis introduced a profile-based resistant-client simulation for questionnaire-driven mental health assessment. The method combines 12 resistance subcategories with turn-level activation, intensity control, prompt boundaries, truth maps for reframing, and metadata that records when resistance is applied. The evaluation covered 1,225 conversations across five clinical profiles, three probabilistic intensity settings, and an extreme stress test.

Resistance reduced question responsiveness and clinically useful information more strongly than the broader conversation-rubric score. The resistance categories also produced different response patterns. Classification accuracy remained above 0.96 from low through high intensity even as diagnostic certainty declined, but it fell sharply in the extreme stress test. Even under this failure condition, the diagnostic agent still produced specific diagnoses, and illustrative inspection showed that an unanswered question could be treated as an endorsed symptom.

Overall, completing the questionnaire did not guarantee that the dialogue contained enough useful information, and a correct diagnostic label did not guarantee that the report was well supported by the patient's answers. Resistant-client simulation made both gaps visible. Future work should validate the generated resistance with human raters and test whether linking diagnostic claims to patient responses reduces unsupported conclusions.

%%% Local Variables:
%%% mode: latex
%%% TeX-master: "../thesis"
%%% End:
